\section{Correctness and Complexity Analysis}
\label{sec:correctness-analysis}

\subsection{Proof Architecture}

The preceding constructions form one proof chain rather than four independent
encodings.  SOLE, in Section~\ref{sec:sole}, exhibits a reversible local change
of finite radix.  Lemma~\ref{lem:information-carrier} isolates the invariant
behind that example: an old spill is committed to an integer number of memory
bits, while a bounded new spill carries only the unresolved part forward.  The
vector result organizes this invariant in an implicit tree, and the online
prefix-free result organizes it in a stream with delayed termination
\cite{dodis2010changing}.  This section audits the proof obligations that are
shared across those applications and records the computational assumptions on
which their time bounds depend.

\subsection{Correctness of the SOLE Transform}

Let \(B=2^b\), let \([B]=\{0,\ldots,B-1\}\), and identify
\(\mathsf{eof}\) with \(B\), so an enlarged-alphabet symbol belongs to
\([B+1]\).  In Pass~1, the pair \((x,y)\in[B+1]^2\) is combined as
\[
    z=y(B+1)+x
\]
and divided by \(B-4i\):
\[
    a=z\bmod(B-4i),
    \qquad
    q=\left\lfloor\frac{z}{B-4i}\right\rfloor.
\]
The remainder automatically satisfies \(a\in[B-4i]\).  The range of the
quotient follows from the full capacity calculation
\[
\begin{aligned}
 &(B-4i)(B+4i+4)-(B+1)^2\\
 &\qquad=2B-16i(i+1)-1.
\end{aligned}
\]
Thus
\[
    (B+1)^2\leq(B-4i)(B+4i+4)
\]
is equivalent to
\[
    B\geq 2(2i+1)^2-\frac{3}{2}.
\]
Basic SOLE assumes \(B\geq2n^2\), and every processed pair satisfies
\(2i+1\leq n\), so the required inequality holds.  Since
\(z<(B+1)^2\), it follows that \(q\in[B+4i+4]\).

The inverse is explicit.  From \((a,q)\), it reconstructs
\[
    z=q(B-4i)+a,
\]
followed by
\[
    x=z\bmod(B+1),
    \qquad
    y=\left\lfloor\frac{z}{B+1}\right\rfloor.
\]
Both uniqueness claims are instances of Euclidean division.  After Pass~1,
Pass~2 regroups a long and a short value.  For \(i\geq1\), the pair belongs to
\([B+4i]\times[B-4i]\), and
\[
    (B+4i)(B-4i)=B^2-16i^2<B^2.
\]
It therefore has a unique representation as two base-\(B\) digits.  The first
\([B]\) value is not part of this shifted pairing, so the strict inequality is
used only where \(i\geq1\).

For the basic odd-length presentation, write \(n=2r+1\).  The appended
\(\mathsf{eof}\) is at position \(2r+2\), and the algorithm appends
\[
    0\in[B-4r-4]
\]
after Pass~1 to complete the final long--short pair.  For even \(n\), the
basic presentation pads to the odd case and can use one more boundary block;
the tighter variant in the paper handles both parities with an early EOF
mechanism.  In either case, termination needs more than the existence of a
distinguished symbol.  SOLE's Property~3 states that input blocks \(2i+1\)
and \(2i+2\) are decoded from output blocks \(2i,\ldots,2i+3\).  The window
that reveals the final EOF lies inside the produced encoding, so the decoder
stops without probing beyond the file.  The same property shows the exact
locality guarantee: encoding and decoding use four-block neighborhoods, not a
single block.

\subsection{Correctness and Redundancy of the Information-Carrier Lemma}

Section~\ref{sec:information-carriers} gives the construction; here we make
its proof obligations explicit.  For \(x\in[X]\), \(y\in[Y]\), choose
\(M\) so that
\[
    C=\left\lfloor\frac{2^M}{Y}\right\rfloor
    =\Theta(\sqrt X),
    \qquad
    S=\left\lceil\frac{X}{C}\right\rceil,
\]
with \(C\geq1\).  Define
\[
    c=x\bmod C,
    \qquad
    s=\left\lfloor\frac{x}{C}\right\rfloor,
    \qquad
    m=yC+c.
\]
Because \(0\leq c<C\) and \(0\leq y<Y\),
\[
    0\leq m<YC\leq2^M,
\]
so \(m\) is a legal \(M\)-bit memory value.  Euclidean division gives
\[
    y=\left\lfloor\frac{m}{C}\right\rfloor,
    \qquad
    c=m\bmod C,
    \qquad
    x=sC+c.
\]
Consequently the map is injective.  More importantly, the old spill \(y\) is
recovered from \(m\) alone; recovering it does not initiate a chain through
later spills.  This asymmetric local recovery condition is what supports the
constant-time accesses in both applications.

The two rounding losses can be bounded without suppressing the floor and
ceiling.  From the definition of \(C\),
\[
    YC\leq2^M<Y(C+1),
\]
and hence
\[
\begin{aligned}
    R_1
    &=M-\log_2(YC)\\
    &=\log_2\frac{2^M}{YC}
      <\log_2\left(1+\frac{1}{C}\right)
      =O\left(\frac{1}{C}\right).
\end{aligned}
\]
Similarly, \(S=\lceil X/C\rceil\) implies \(CS<X+C\), so
\[
\begin{aligned}
    R_2
    &=\log_2C+\log_2S-\log_2X\\
    &=\log_2\frac{CS}{X}
      <\log_2\left(1+\frac{C}{X}\right)
      =O\left(\frac{C}{X}\right).
\end{aligned}
\]
Balancing the two terms at \(C=\Theta(\sqrt X)\) yields
\[
    R_1+R_2=O\left(\frac{1}{\sqrt X}\right),
    \qquad
    S=O(\sqrt X),
    \qquad
    2^M=O(Y\sqrt X).
\]
The encoder and decoder use a constant number of word additions,
multiplications, divisions, and remainder operations once \(M,C,S\) are
precomputed.

\subsection{Proof of the Optimal-Space Vector Result}

In Section~\ref{sec:vector-representation}, \(\Sigma\) denotes the alphabet
cardinality.  Consecutive coordinates are grouped into carriers from a
universe
\[
    X=\Theta(n^2).
\]
For a fixed finite alphabet of size \(\Sigma\geq2\), each carrier contains
\(\Theta(\log_{\Sigma}n)\) coordinates and the number of carriers is
\[
    N=O\left(\frac{n}{\log_{\Sigma}n}\right).
\]
Lemma~\ref{lem:information-carrier} contributes \(O(1/n)\) redundancy per
carrier, so the total carrier redundancy is
\[
    N\,O\left(\frac{1}{n}\right)
    =O\left(\frac{1}{\log_{\Sigma}n}\right)
    =o(1).
\]

The sequential first attempt already has this space bound and local access.
Its failure is nonuniform precomputation: irregular spill universes can
require \(O(N)\) distinct parameter sets and memory offsets.  The implicit
binary tree reduces that cost.  On any level its nodes form at most three
classes: complete subtrees of height \(k\), at most one incomplete subtree,
and complete subtrees of height \(k-1\).  Nodes in one class receive identical
child-spill universe sizes and therefore share carrier parameters.  There are
only constantly many classes per level and \(O(\log n)\) levels, giving the
\(O(\log n)\) precomputed word constants in
Theorem~\ref{thm:vector-representation}.

The tree does not cause logarithmic query time.  To recover the carrier at a
nonroot node \(v\), the algorithm uses the parent's memory to recover the
spill sent by \(v\), then combines that spill with \(v\)'s own memory.  It
touches the node and its parent, not a root-to-leaf path.  The root uses its
separately stored spill.  A coordinate is then extracted from the recovered
carrier by fixed-radix word arithmetic.

An update is equally local.  Re-encoding \(v\) may change its outgoing spill,
so its parent must also be re-encoded.  In the carrier construction, however,
the new spill is the carrier quotient
\[
    s=\left\lfloor\frac{x}{C}\right\rfloor,
\]
which does not depend on the old-spill value.  The parent's own carrier and
parameters are unchanged; therefore its outgoing spill is unchanged even
though its memory content absorbs the modified child spill.  No update reaches
the grandparent.  Roots, leaves, and the unique incomplete subtree use the
same argument with separately stored or singleton missing-child spills.

Finally, \(n\log_2\Sigma+o(1)\) is not by itself the claimed exact integer
bound.  Storing the final root spill gives initially
\[
    \left\lceil n\log_2\Sigma+o(1)\right\rceil,
\]
which could be one bit larger than \(\lceil n\log_2\Sigma\rceil\).  The paper
enlarges the carriers so that the residual redundancy is \(O(n^{-c})\) for
any fixed \(c\), while arithmetic still occupies a constant number of
\(O(\log n)\)-bit words.  For fixed \(\Sigma\) that is not a power of two, a
linear-forms-in-logarithms bound gives an inverse-polynomial lower bound on
the nonzero distance from \(n\log_2\Sigma\) to an integer.  Choosing \(c\)
large enough keeps the residual below the next-integer gap, so the final
ceiling absorbs it.  Power-of-two alphabets have an integral information
bound and admit direct fixed-width packing.  The number-theoretic theorem is
used only for this last rounding step, not for locality or the carrier lemma.

\subsection{Proof of the Online Prefix-Free Result}

To avoid overloading \(n\), let \(t\) be the number of bits observed when the
next block size is chosen, let \(n\) be the final input length, and let \(N\)
be the final number of blocks.  A complete block has two halves of length
\[
    k_i=\lceil\log_2t\rceil
\]
and hence universe \(X_i=2^{2k_i}=\Theta(t^2)\), up to scale constants.  The
carrier loss is therefore \(O(2^{-k_i})\).  Enlarging one half-alphabet from
\([2^{k_i}]\) to \([2^{k_i}]\cup\{\mathsf{eof}\}\) costs
\[
\begin{aligned}
 &\log_2\bigl(2^{k_i}(2^{k_i}+1)\bigr)-2k_i\\
 &\qquad=\log_2(1+2^{-k_i})=O(2^{-k_i}).
\end{aligned}
\]
At scale \(k_i=K\), there are \(O(2^K/K)\) blocks because the scale contains
\(O(2^K)\) input bits and each block has \(\Theta(K)\) bits.  Each of the two
losses thus contributes \(O(1/K)\) at that scale.  Since the largest scale is
\(O(\log n)\),
\[
    \sum_{K=1}^{\lceil\log_2n\rceil}O\left(\frac{1}{K}\right)
    =O(\log\log n).
\]

Prefix-freeness depends on the timing of EOF discovery.  A three-block lag
keeps a constant number of blocks in an \(O(\log n)\)-bit buffer.  If complete
input block \(N-2\) was \((a,b)\), the encoder replaces it by
\((a,\mathsf{eof})\) and saves the displaced half \(b\) in the special tail.
The last incomplete block \(N\) is also encoded in that tail with an Elias
code.  No already emitted output is rewritten.

After reading ordinary output block \(N-1\), the sequential carrier inverse
has the spill required to reconstruct input block \(N-2\).  It therefore sees
\(\mathsf{eof}\) before requesting another nonexistent ordinary output block,
stops the carrier decoder, and switches to the self-delimiting tail.  The tail
restores \(b\) and the final partial block.  The length decomposition is
\[
\begin{aligned}
    \text{length}
    &=\underbrace{n}_{\text{payload}}
      +\underbrace{\log_2n}_{\text{displaced half / early EOF}}\\
    &\quad+\underbrace{O(\log\log n)}_{\text{alphabet enlargement}}
      +\underbrace{O(\log\log n)}_{\text{carrier rounding}}\\
    &\quad+\underbrace{O(\log\log n)}_{\text{incomplete-block tail}}\\
    &=n+\log_2n+O(\log\log n).
\end{aligned}
\]
The displaced half and the early EOF are one mechanism and are not charged as
two separate \(\log n\) terms.  This is near-optimal, within
\(O(\log\log n)\) of the Elias-scale bound; it is not an exact reproduction
of the full Elias omega iterated-log expression.

The streaming state consists of a current block, a bounded spill, block-size
counters, the constant-lag buffer, and termination metadata.  These are a
constant number of \(O(\log n)\)-bit objects, so the working space is
\(O(\log n)\) bits.  The time guarantee is \(O(1)\) per word-level or block
operation, not per individual bit.

\subsection{Computational Model and Complexity Assumptions}

All constant-time claims use the paper's Word RAM model.  Memory supports
random access and is divided into \(w\)-bit words with
\(w=\Omega(\log n)\), so indices and the ordinary carrier parameters fit in
a word.  Addition, multiplication, and division are unit-cost operations;
remainder is obtained with quotient--remainder arithmetic.  Each local
transformation uses a constant number of words, and its universe-dependent
constants are precomputed.  The vector data structure counts
\(O(\log n)\) such constants explicitly.  The exact-space refinement uses
\(O(c\log n)\)-bit integers for fixed \(c\), still a constant number of
words in the asymptotic model.

These are theoretical model guarantees.  They do not imply that hardware
division, cache behavior, or a particular implementation has been benchmarked
to be faster on every real platform.  Likewise, the online bound is stated
per block or word operation; translating it to a per-bit engineering cost
depends on the chosen block schedule and machine word size.

\subsection{Scope of the Guarantees}

The paper proves an exact-space representation for the stated uniform
finite-alphabet vector, constant-time coordinate reads and writes on its Word
RAM, a low-memory online prefix-free code, and local reversible base
conversion.  It does not prove a locally decodable arithmetic code for biased
distributions, a zero-redundancy rank/select structure, or a general
dictionary solution.  Nor does prefix-free preprocessing automatically prove
the security of every cryptographic construction: the primitive, padding,
and security model still require a separate argument, as discussed in
Section~\ref{sec:related-work}.  The results also do not assert optimality for
arbitrary compression models or superior performance on every physical
machine.
